%! Logarithmic Functions — Unit 2, Lesson 2.11 — Group Activity
\documentclass[10pt]{article}
\usepackage{precalculus-article}
\usepackage{precalculus-boxes}
\newcommand{\TallMath}[1]{$\displaystyle #1\rule[-1.4em]{0pt}{3.2em}$}

\begin{document}

\pageheader{Unit 2, Lesson 2.11}{Group Activity}
\namepartnerperiod

\vspace{0.08in}

\begin{scenariobox}[Seismic Data Analysis]{sky}
A geologist uses $f(x) = \log_{10}(x)$ to interpret seismic intensity readings, where $x$ is
the measured wave intensity. Your group will analyze the characteristics of logarithmic functions
in this context and beyond.
\end{scenariobox}

\vspace{0.08in}

\begin{tcolorbox}[colback=white, colframe=black!40,
    title=\textbf{Tier R --- Remediate},
    fonttitle=\bfseries, arc=2mm, left=3mm, right=3mm, top=2mm, bottom=2mm]

\textbf{Work with $f(x) = \log_{10}(x)$.}

\begin{enumerate}[leftmargin=1.8em, itemsep=10pt]

\item State the domain and range of $f(x) = \log_{10}(x)$.

\vspace{4pt}
Domain: \blank{4cm} \qquad Range: \blank{4cm}

\vspace{4pt}
Equation of the vertical asymptote: $x = $ \blank{2cm}

\item Complete the table of values.

\vspace{4pt}
\begin{center}
\renewcommand{\arraystretch}{1.8}
\begin{tabular}{|c|c|}
\hline
\rowcolor{sky}
$x$ & $f(x) = \log_{10}(x)$ \\
\hline
$0.1$   & \blank{2.5cm} \\
\hline
$1$     & \blank{2.5cm} \\
\hline
$10$    & \blank{2.5cm} \\
\hline
$100$   & \blank{2.5cm} \\
\hline
$1000$  & \blank{2.5cm} \\
\hline
\end{tabular}
\end{center}

\item The geologist records the readings $x = -5$, $x = 0$, and $x = 7.3$.
      Which readings are \textbf{invalid} inputs for $f$? Explain why.

\writelines{2}

\item Based on your table, is $f$ increasing or decreasing on its domain?

\vspace{4pt}
$f$ is \blank{3cm} on $(0, \infty)$.

\end{enumerate}
\end{tcolorbox}

\vspace{0.08in}

\begin{tcolorbox}[colback=white, colframe=black!40,
    title=\textbf{Tier A --- Approaching Proficiency},
    fonttitle=\bfseries, arc=2mm, left=3mm, right=3mm, top=2mm, bottom=2mm]

\textbf{Work with $f(x) = 2\log_3(x)$.}

\begin{enumerate}[leftmargin=1.8em, itemsep=10pt]

\item State the domain and range. Write the equation of the vertical asymptote.

\vspace{4pt}
Domain: \blank{3.5cm} \quad Range: \blank{3.5cm} \quad Asymptote: $x = $ \blank{1.5cm}

\item Fill in the end behavior using limit notation.

\vspace{4pt}
$\displaystyle\lim_{x \to 0^+} 2\log_3(x) = $ \blank{2.5cm} \qquad
$\displaystyle\lim_{x \to \infty} 2\log_3(x) = $ \blank{2.5cm}

\item Is $f$ increasing or decreasing? How do you know from the inverse relationship?

\writelines{2}

\item Is the graph of $f$ concave up or concave down? Does $f$ have any local maximum or
      minimum on its natural domain? Explain using EK 2.11.A.2.

\writelines{2}

\item Explain in words what happens to the output of $f(x) = 2\log_3(x)$ as $x$ gets
      very close to zero from the right.

\writelines{2}

\end{enumerate}
\end{tcolorbox}

\vspace{0.08in}

\begin{tcolorbox}[colback=white, colframe=black!40,
    title=\textbf{Tier E --- Extension},
    fonttitle=\bfseries, arc=2mm, left=3mm, right=3mm, top=2mm, bottom=2mm]

\textbf{Additive Transformation Identification Test.}

\begin{enumerate}[leftmargin=1.8em, itemsep=10pt]

\item The table below shows $f(x) = \log_3 x$. Check whether the inputs are proportional
      over equal-length output intervals by computing each ratio.

\vspace{4pt}
\begin{center}
\renewcommand{\arraystretch}{1.8}
\begin{tabular}{|c|c|c|}
\hline
\rowcolor{sky}
$x$ & $f(x)$ & Input ratio (consecutive) \\
\hline
$1$  & $0$ & --- \\
\hline
$3$  & $1$ & \blank{3cm} \\
\hline
$9$  & $2$ & \blank{3cm} \\
\hline
$27$ & $3$ & \blank{3cm} \\
\hline
$81$ & $4$ & \blank{3cm} \\
\hline
\end{tabular}
\end{center}

\vspace{4pt}
Are the inputs proportional? \blank{2cm} \quad Constant ratio: \blank{2cm}

\item Now form $g(x) = f(x + 2) = \log_3(x + 2)$. The output values are still $0,1,2,3,4$,
      but the inputs that produce those outputs change. Find the new inputs by solving
      $\log_3(x + 2) = k$ for each output $k = 0, 1, 2, 3, 4$.

\vspace{4pt}
\begin{center}
\renewcommand{\arraystretch}{1.8}
\begin{tabular}{|c|c|c|}
\hline
\rowcolor{sky}
$x$ (new input for $g$) & $g(x)$ & Input ratio (consecutive) \\
\hline
\blank{2.5cm} & $0$ & --- \\
\hline
\blank{2.5cm} & $1$ & \blank{3cm} \\
\hline
\blank{2.5cm} & $2$ & \blank{3cm} \\
\hline
\blank{2.5cm} & $3$ & \blank{3cm} \\
\hline
\blank{2.5cm} & $4$ & \blank{3cm} \\
\hline
\end{tabular}
\end{center}

\vspace{4pt}
Are the inputs of $g$ proportional? \blank{2cm} \quad This confirms EK 2.11.A.3.

\item A mystery function $h$ has the following table. Determine whether $h$ could be
      logarithmic by testing whether the inputs are proportional over equal output intervals.

\vspace{4pt}
\begin{center}
\renewcommand{\arraystretch}{1.8}
\begin{tabular}{|c|c|}
\hline
\rowcolor{sky}
$x$ & $h(x)$ \\
\hline
$2$  & $1$ \\
\hline
$8$  & $2$ \\
\hline
$32$ & $3$ \\
\hline
$128$& $4$ \\
\hline
\end{tabular}
\end{center}

\vspace{4pt}
Input ratios: \blank{5cm}

\vspace{4pt}
Is $h$ logarithmic? Justify. \writelines{2}

\end{enumerate}
\end{tcolorbox}

\end{document}
